% ----------------------------------------------------------
% Introdução
% ----------------------------------------------------------
\chapter{INTRODUÇÃO}

Em 2022, a indústria de jogos eletrônicos movimentou cerca de US\$ 197 bilhões, superando a indústria do cinema e ocupando o segundo lugar nos negócios de entretenimento mundial, atrás apenas da TV \cite{costa2024evolucao}, além de reunir uma base de jogadores que já ultrapassa a marca de três bilhões de pessoas ao redor do mundo \cite{newzoo2025}. Esse crescimento tem transformado tanto o perfil do consumidor quanto a forma como os jogos são pensados e desenvolvidos, usando modelos mais colaborativos de produção.

Nesse contexto, o \textit{feedback} da comunidade tem se tornado essencial no ciclo de desenvolvimento de jogos digitais, especialmente para desenvolvedores independentes, que conseguem nas contribuições dos jogadores uma fonte acessível de informações para guiar suas decisões criativas e técnicas ao longo das fases do projeto desde a concepção até a pós-produção \cite{uchoa2018}. Essa prática se popularizou com a metodologia de Desenvolvimento Aberto (\textit{Open Development}), na qual desenvolvedores disponibilizam publicamente versões incompletas dos jogos e as aprimoram continuamente com base no retorno da comunidade \cite{thominet2018}.

A importância desse modelo é ressaltada por \citeonline{vargas2018} que, ao analisarem o caso do \textit{Pokémon Go}, demonstraram que a Niantic aderiu às demandas recorrentes da comunidade em atualizações do jogo, indicando que o \textit{feedback} dos jogadores, quando coletado e analisado, pode se converter em evolução para o produto. \citeonline{moro2022} também reforçam que a comunidade on-line é capaz de gerar uma quantidade de \textit{feedbacks} superior à obtida em ambientes mais restritos.

Entretanto, a indústria enfrenta um desafio para o aproveitamento desse potencial: a dispersão entre diversas plataformas e a variação de perfis existentes nelas. \citeonline{tong2022} mostra que usuários da plataforma Steam tendem a focar em aspectos técnicos e mecânicas de jogo, enquanto, no YouTube, o foco costuma recair sobre elementos estéticos, visuais e narrativos. Essa fragmentação dificulta que os desenvolvedores analisem e compreendam as demandas da comunidade de maneira eficiente.

Como resultado prático desta pesquisa, será desenvolvido o protótipo de uma plataforma web denominada Spaw Ideias, destinada à interação entre desenvolvedores e jogadores de diversos segmentos para o compartilhamento de ideias, sugestões e feedbacks sobre jogos. O sistema também se propõe a reunir dados que ofereçam informações valiosas para que empresas e desenvolvedores de jogos individuais possam tomar decisões informadas sobre as características e o escopo de seus produtos.
